\chapter*{Jak czytać wspólny skrypt}
\addcontentsline{toc}{chapter}{Jak czytać wspólny skrypt}
Najpierw wprowadzamy grafy, związane z nimi macierze i realizacje topologiczne.
Następnie przechodzimy do geometrii rozmaitości, działań grup,
węzłów, chirurgii i topologii niskowymiarowej.
Wszystkie diagramy są edytowalne w TikZ. Dowody możliwe do przeprowadzenia
przy pomocy poznanego aparatu rozwijamy bez ukrywania istotnych kroków;
twierdzenia wymagające obszernej teorii dodatkowej oznaczamy jako
wyniki przyjęte z literatury. Przy rezultatach korzystających z teorii
indeksu, zwartości lub transwersalności rozróżniamy rachunek formalny,
skonstruowanie obiektu i zastosowanie zewnętrznego twierdzenia.
W całym tekście używamy polskiego terminu ``zanurzenie'' w znaczeniu
\emph{embedding}, odróżniając go od immersji.

